\documentclass[11pt,a4paper]{article}
\usepackage[T1]{fontenc}
\usepackage[utf8]{inputenc}
\usepackage{amsmath,amssymb,amsthm}
\usepackage[margin=1in]{geometry}
\usepackage[colorlinks=true,linkcolor=blue,urlcolor=blue]{hyperref}
\theoremstyle{plain}
\newtheorem{theorem}{Theorem}
\newtheorem{lemma}[theorem]{Lemma}
\newtheorem{proposition}[theorem]{Proposition}
\newtheorem{corollary}[theorem]{Corollary}
\theoremstyle{definition}
\newtheorem{remark}[theorem]{Remark}

\title{Recovering the unwritten recurrences of the table OEIS A234343}
\author{Adrian Perez Fontelles\\ \small Independent researcher}
\date{2 September 2026}
\begin{document}
\maketitle

\begin{abstract}
OEIS A234343 is a two-dimensional table $T(n,k)$ whose empirical recurrences are, for several of
its lines, given only as an ORDER --- ``k=2: [order 77]'' and the like --- with the
coefficients in a linked file rather than in the entry. Those conjectures can still be settled,
and the recurrences recovered. A line of the table is a fixed-width or fixed-height array
count, hence a walk count on finitely many vertices, hence satisfies some monic recurrence of
order at most that number; Berlekamp--Massey on twice as many exact terms returns the MINIMAL
one. Where its order equals the order the entry states --- and where the same holds after the
first few terms are dropped --- any recurrence of that order the line satisfies has a
characteristic polynomial that is a multiple of the minimal one and of equal degree, hence is
that polynomial. This note settles column 2.
\end{abstract}

\noindent\small 2020 Mathematics Subject Classification. 05A15, 11B37, 15A18.\normalsize

\section{The table and the unwritten conjectures}

OEIS A234343 is ``T(n,k)=Number of (n+1)X(k+1) 0..7 arrays with every 2X2 subblock having the sum of the absolute values of the edge differences equal to 14.''. It is read by antidiagonals and begins
\[
604, 9076, 9076, 120838, 411064, 120838, 1709528, 15724414,\ \dots
\]
The entry carries, among its empirical claims:
\begin{quote}\small
k=2: [order 77]
\end{quote}
As of the ``Last modified'' line on the live entry (Jun 02 2025 09:17:21, revision 6) these are still
recorded as empirical, and the recurrences themselves are not in the entry text.

\section{Why a line of the table is a walk count}

Fix the index that the line fixes. The entry defines $T(n,k)$ as a number of arrays of a shape
determined by $n$ and $k$, subject to a condition local to a bounded window of consecutive
lines. Substituting the fixed index into the entry's own wording turns the definition into an
ordinary one-parameter array count, and for such a count the admissible windows are the
vertices of a finite digraph and an array is a walk. So the line is a walk count on $S$
vertices, and by the Cayley--Hamilton theorem it satisfies a monic integer recurrence of order
at most $S$.

\section{Recovering the recurrences}

\begin{lemma}\label{lem:bm}
A sequence known to satisfy some linear recurrence of order at most $S$ is determined, with its
minimal recurrence, by its first $2S$ terms; Berlekamp--Massey applied to them returns that
minimal recurrence exactly.
\end{lemma}

\subsection*{Column $k=2$}
Substituting the fixed index into the entry's wording gives an ordinary one-parameter array
count, a walk on $S=512$ vertices. Berlekamp--Massey on $2S=1024$ exact terms
returns a minimal recurrence of order $77$ --- the order the entry states --- with
integer coefficients
\[
(c_1,\dots,c_{77})=(46,\ 686,\ -39131,\ -147608,\ 13161758,\ 5491556,\ -2422071738,\ 2669301671,\ 280934675228,\ -544179239981,\ -22251655048299,\ 54720382635220,\ 1267432845547909,\ -3576653091034131,\ -53817823438579862,\ 166796201636618054,\ 1748506550387989431,\ -5824255814758095441,\ -44303157179923285861,\ 156893139869964306677,\ 887710616704701640930,\ -3326540630701289362292,\ -14205591683763300259754,\ 56297099545629304041237,\ 182735902915047707420017,\ -767978059616287745836565,\ -1896332022514815683832791,\ 8501727703353782268816913,\ 15888408497295948336347226,\ -76708964216142481329364980,\ -107260449257105303974342660,\ 565472654093678652002345121,\ 580358573439870006956132749,\ -3408463577123319957190881689,\ -2492850609278114242335730243,\ 16790757499970745612792206053,\ 8363741285815559104050258937,\ -67493887881664764986109749144,\ -21296199132742063284311047699,\ 220851309477824621103121309649,\ 38761362519510759680396864617,\ -586525465202239901670293792807,\ -42211587509722646041687816792,\ 1260075442354729229942951363119,\ 166592726454454084372210872,\ -2182500915576822076327099787940,\ 96071521254690064489893640134,\ 3037262857751753136498131019873,\ -194978093453624708331263307792,\ -3384177795462660213307834500786,\ 217011288989783962536339115686,\ 3007012515559713352976779366246,\ -143474728321082619553140285259,\ -2120011712535679103062495702038,\ 40296428707275175849286144706,\ 1177861084554219483602981121684,\ 21208736542570154113267479073,\ -510855824006477751691754030283,\ -30505073886353428786157245725,\ 170722777333502275178118799597,\ 17616871680448863839036577568,\ -43179766888262456598409730485,\ -6324085849369424237145168690,\ 8060920992777078722215966477,\ 1521329718161266593417803797,\ -1071104515056525255967624652,\ -247228926969628374230067392,\ 95721024748194936713263392,\ 26474131406274123724436256,\ -5195532295082729446156496,\ -1761977505868088586782800,\ 132916098379404803890848,\ 64761631889591477372608,\ 175828816058823338496,\ -979982859324325789440,\ -51737541426085860864,\ -37734968676925440).
\]
so that $a(m)=\sum_{i=1}^{77}c_i\,a(m-i)$ for every $m$. Removing the first
$77+4$ terms and repeating gives order $77$ again, which is what the
identification below needs.

\begin{theorem}
For each line above, the displayed recurrence holds for every index, and it is the recurrence
the entry records.
\end{theorem}

\begin{proof}
It holds by Lemma~\ref{lem:bm}. For the identification, the entry asserts that the line
satisfies SOME recurrence of the stated order, possibly only from some index on. Let $q$ be its
characteristic polynomial and $q_{\min}$ the minimal polynomial of the tail. Then
$q_{\min}\mid q$, and the second Berlekamp--Massey run --- on the line with its first few
terms removed --- shows $\deg q_{\min}$ equals the stated order, which is $\deg q$. Both are
monic, so $q=q_{\min}$.
\end{proof}

\section{Verification}

Three checks.

First, each line's model was compared against the entry's own published values for that line,
taken from the antidiagonal DATA read in either orientation and from the ``Table starts''
block, and agrees at every available term. This is the only point at which reading the entry's
English enters, and a misreading gives different counts.

Second, each recovered recurrence was evaluated directly on the model's values, with no
Berlekamp--Massey involved, and holds at every index tested.

Third, each order was computed twice by independent routes --- modulo a $61$-bit prime, where
every intermediate is a machine word, and again in exact rational arithmetic --- and once more
on the line with its first few terms dropped, which is what the identification requires. All
agree.

\begin{thebibliography}{9}
\bibitem{oeis} The OEIS Foundation, \emph{The On-Line Encyclopedia of Integer
Sequences}, \url{https://oeis.org}, 2026. Sequence A234343.
\bibitem{massey} J.~L.~Massey, Shift-register synthesis and BCH decoding, \emph{IEEE
Trans. Inform. Theory} \textbf{15} (1969), 122--127.
\bibitem{stanley} R.~P.~Stanley, \emph{Enumerative Combinatorics}, Volume 1, 2nd ed.,
Cambridge University Press, 2011. (Transfer-matrix method, Section 4.7.)
\bibitem{horn} R.~A.~Horn and C.~R.~Johnson, \emph{Matrix Analysis}, 2nd ed., Cambridge
University Press, 2013.
\end{thebibliography}

\end{document}
